\documentclass[11pt,a4paper]{article}
\usepackage[T1]{fontenc}
\usepackage[utf8]{inputenc}
\usepackage{amsmath,amssymb,amsthm}
\usepackage[margin=1in]{geometry}
\usepackage[colorlinks=true,linkcolor=blue,urlcolor=blue]{hyperref}
\theoremstyle{plain}
\newtheorem{theorem}{Theorem}
\newtheorem{lemma}[theorem]{Lemma}
\newtheorem{proposition}[theorem]{Proposition}
\newtheorem{corollary}[theorem]{Corollary}
\theoremstyle{definition}
\newtheorem{remark}[theorem]{Remark}

\title{Recovering the unwritten recurrences of the table OEIS A236798}
\author{Adrian Perez Fontelles\\ \small Independent researcher}
\date{2 September 2026}
\begin{document}
\maketitle

\begin{abstract}
OEIS A236798 is a two-dimensional table $T(n,k)$ whose empirical recurrences are, for several of
its lines, given only as an ORDER --- ``k=2: [order 38]'' and the like --- with the
coefficients in a linked file rather than in the entry. Those conjectures can still be settled,
and the recurrences recovered. A line of the table is a fixed-width or fixed-height array
count, hence a walk count on finitely many vertices, hence satisfies some monic recurrence of
order at most that number; Berlekamp--Massey on twice as many exact terms returns the MINIMAL
one. Where its order equals the order the entry states --- and where the same holds after the
first few terms are dropped --- any recurrence of that order the line satisfies has a
characteristic polynomial that is a multiple of the minimal one and of equal degree, hence is
that polynomial. This note settles column 2, column 3.
\end{abstract}

\noindent\small 2020 Mathematics Subject Classification. 05A15, 11B37, 15A18.\normalsize

\section{The table and the unwritten conjectures}

OEIS A236798 is ``T(n,k)=Number of (n+1)X(k+1) 0..2 arrays with the maximum plus the upper median plus the minimum of every 2X2 subblock equal.''. It is read by antidiagonals and begins
\[
81, 255, 255, 825, 1231, 825, 2817, 6217,\ \dots
\]
The entry carries, among its empirical claims:
\begin{quote}\small
k=2: [order 38]\\
k=3: [order 83]
\end{quote}
As of the ``Last modified'' line on the live entry (Jul 23 2025 09:33:56, revision 6) these are still
recorded as empirical, and the recurrences themselves are not in the entry text.

\section{Why a line of the table is a walk count}

Fix the index that the line fixes. The entry defines $T(n,k)$ as a number of arrays of a shape
determined by $n$ and $k$, subject to a condition local to a bounded window of consecutive
lines. Substituting the fixed index into the entry's own wording turns the definition into an
ordinary one-parameter array count, and for such a count the admissible windows are the
vertices of a finite digraph and an array is a walk. So the line is a walk count on $S$
vertices, and by the Cayley--Hamilton theorem it satisfies a monic integer recurrence of order
at most $S$.

\section{Recovering the recurrences}

\begin{lemma}\label{lem:bm}
A sequence known to satisfy some linear recurrence of order at most $S$ is determined, with its
minimal recurrence, by its first $2S$ terms; Berlekamp--Massey applied to them returns that
minimal recurrence exactly.
\end{lemma}

\subsection*{Column $k=2$}
Substituting the fixed index into the entry's wording gives an ordinary one-parameter array
count, a walk on $S=78$ vertices. Berlekamp--Massey on $2S=156$ exact terms
returns a minimal recurrence of order $38$ --- the order the entry states --- with
integer coefficients
\[
(c_1,\dots,c_{38})=(15,\ -9,\ -880,\ 3426,\ 18155,\ -121505,\ -115924,\ 2003052,\ -1313613,\ -18570720,\ 30811261,\ 100624740,\ -269758235,\ -288602304,\ 1357891776,\ 124420915,\ -4267235783,\ 2172222593,\ 8388777279,\ -8441776674,\ -9593628558,\ 15912829245,\ 4511594574,\ -17273359586,\ 2518712071,\ 10814653852,\ -4621110462,\ -3640874304,\ 2541569180,\ 539526448,\ -683839304,\ 5712224,\ 95082784,\ -11026368,\ -6334080,\ 1095936,\ 148480,\ -29696).
\]
so that $a(m)=\sum_{i=1}^{38}c_i\,a(m-i)$ for every $m$. Removing the first
$38+4$ terms and repeating gives order $38$ again, which is what the
identification below needs.

\subsection*{Column $k=3$}
Substituting the fixed index into the entry's wording gives an ordinary one-parameter array
count, a walk on $S=190$ vertices. Berlekamp--Massey on $2S=380$ exact terms
returns a minimal recurrence of order $83$ --- the order the entry states --- with
integer coefficients
\[
(c_1,\dots,c_{83})=(17,\ 149,\ -3822,\ -5099,\ 385130,\ -470186,\ -23210949,\ 60570417,\ 940504129,\ -3379677366,\ -27246699520,\ 121431917218,\ 584194818792,\ -3139393787896,\ -9425097731306,\ 61598162124697,\ 114219500453427,\ -946862834116793,\ -1004826003316698,\ 11646036448147900,\ 5571787380506985,\ -116340952161916838,\ -3272779809783538,\ 954302011875230349,\ -320548644721837094,\ -6478914470291856511,\ 4170768349052004296,\ 36612702131063326183,\ -33574882254880366927,\ -172831891596991986931,\ 202795687958598677053,\ 682536774766893321784,\ -975834084459192238409,\ -2252938899441608549659,\ 3843391626158026393232,\ 6190401975549878168591,\ -12570611244671849135022,\ -14030566709588819286664,\ 34427622778574228358699,\ 25747525903434993537460,\ -79324604036960951397983,\ -36717985022501057720706,\ 154113240952271945038052,\ 36243384479322799503009,\ -252558759358735794910697,\ -12090990432014794852664,\ 348668526245119369378700,\ -38616999797218798816301,\ -404308939215042828075753,\ 100905470905950208406915,\ 391915280100857250940160,\ -147383978379028506438470,\ -315346851989417002084179,\ 156919130648950916587258,\ 208477619911831767937308,\ -129927927510832504586596,\ -111524272768219126399088,\ 85511507912413637949661,\ 47101637328217152670628,\ -44993104485233133278692,\ -15004974889836196172242,\ 18870811452105240611561,\ 3225918607639655186313,\ -6248749870809523853017,\ -270423641206479033614,\ 1606626900177725739510,\ -100322171731125546752,\ -312327302479984987304,\ 48944503662932084704,\ 43923239357947742776,\ -11082174452615009040,\ -4104648227612050704,\ 1549756930847606688,\ 201467058895942272,\ -135201595488903808,\ 1500409585375488,\ 6675049646263808,\ -747982718679040,\ -133709763618816,\ 30756323024896,\ -785080713216,\ -215328096256,\ 14144765952).
\]
so that $a(m)=\sum_{i=1}^{83}c_i\,a(m-i)$ for every $m$. Removing the first
$83+4$ terms and repeating gives order $83$ again, which is what the
identification below needs.

\begin{theorem}
For each line above, the displayed recurrence holds for every index, and it is the recurrence
the entry records.
\end{theorem}

\begin{proof}
It holds by Lemma~\ref{lem:bm}. For the identification, the entry asserts that the line
satisfies SOME recurrence of the stated order, possibly only from some index on. Let $q$ be its
characteristic polynomial and $q_{\min}$ the minimal polynomial of the tail. Then
$q_{\min}\mid q$, and the second Berlekamp--Massey run --- on the line with its first few
terms removed --- shows $\deg q_{\min}$ equals the stated order, which is $\deg q$. Both are
monic, so $q=q_{\min}$.
\end{proof}

\section{Verification}

Three checks.

First, each line's model was compared against the entry's own published values for that line,
taken from the antidiagonal DATA read in either orientation and from the ``Table starts''
block, and agrees at every available term. This is the only point at which reading the entry's
English enters, and a misreading gives different counts.

Second, each recovered recurrence was evaluated directly on the model's values, with no
Berlekamp--Massey involved, and holds at every index tested.

Third, each order was computed twice by independent routes --- modulo a $61$-bit prime, where
every intermediate is a machine word, and again in exact rational arithmetic --- and once more
on the line with its first few terms dropped, which is what the identification requires. All
agree.

\begin{thebibliography}{9}
\bibitem{oeis} The OEIS Foundation, \emph{The On-Line Encyclopedia of Integer
Sequences}, \url{https://oeis.org}, 2026. Sequence A236798.
\bibitem{massey} J.~L.~Massey, Shift-register synthesis and BCH decoding, \emph{IEEE
Trans. Inform. Theory} \textbf{15} (1969), 122--127.
\bibitem{stanley} R.~P.~Stanley, \emph{Enumerative Combinatorics}, Volume 1, 2nd ed.,
Cambridge University Press, 2011. (Transfer-matrix method, Section 4.7.)
\bibitem{horn} R.~A.~Horn and C.~R.~Johnson, \emph{Matrix Analysis}, 2nd ed., Cambridge
University Press, 2013.
\end{thebibliography}

\end{document}
